% =====================================================================
\section{Fase de Estabilización}
\label{sec:estabilizacion}

En esta fase no se agregó funcionalidad nueva. El trabajo se concentró en
integrar los módulos, uniformar la arquitectura y verificar la calidad del
código.

\subsection{Integración de los módulos}

Los módulos desarrollados en las ramas \archivo{feature/auth} y
\archivo{feature/machines} se integraron en \archivo{develop} mediante
solicitudes de integración. La navegación principal quedó unificada en la
barra de pestañas de \archivo{(app)} (Inicio, Máquinas, Progreso y
Nutrición), con las rutinas y el perfil accesibles desde el inicio, y la
protección de rutas centralizada en el diseño raíz.

\subsection{Refactorización por capas}

Durante las primeras iteraciones, solo el módulo de nutrición cumplía
completamente la separación en capas; las pantallas de autenticación y de
máquinas llamaban a Firebase directamente. En la estabilización se extendió el
patrón a todo el sistema, de modo que cada módulo cuenta con su controlador,
sus entidades, su esquema de validación y su repositorio
(Tabla~\ref{tab:modulos-capas}).

\begin{table}[H]
\centering
\caption{Componentes de cada módulo por capa}
\label{tab:modulos-capas}
\small
\begin{tabularx}{\textwidth}{|L{2.1cm}|L{4.0cm}|Y|L{3.0cm}|}
\hline
\rowcolor{encabezado} \textbf{Módulo} & \textbf{Controlador} & \textbf{Modelo (entidad /
repositorio)} & \textbf{Vista} \\
\hline
Autenticación & \archivo{AuthController} & \archivo{Usuario} /
\archivo{usuarioRepository} & \archivo{(auth)}, \archivo{(onboarding)},
\archivo{perfil} \\ \hline
Máquinas & \archivo{MaquinasController} & \archivo{Maquina} /
\archivo{maquinaRepository} & \archivo{machines/*} \\ \hline
Ejercicios & \archivo{EjerciciosWgerController} &
\archivo{EjercicioWger} / \archivo{ejercicioWgerRepository} & Modales de
selección \\ \hline
Rutinas & \archivo{RutinasController} & \archivo{Rutina} /
\archivo{rutinaRepository} & \archivo{routines/*} \\ \hline
Nutrición & \archivo{NutritionController} & \archivo{RegistroNutricional}
/ \archivo{nutricionRepository} & \archivo{nutricion} \\ \hline
Progreso & \archivo{ProgresoController} & \archivo{ProgresoFisico} /
\archivo{progresoRepository} & \archivo{progress/*}, inicio \\
\hline
\end{tabularx}
\end{table}

Otras mejoras aplicadas en la refactorización fueron:

\begin{itemize}
    \item \textbf{Tiempo real en lugar de recarga manual.} Las pantallas que
    consultaban los datos una sola vez y ofrecían una función de recarga
    pasaron a suscribirse con \archivo{onSnapshot} a través de los métodos
    \archivo{observar} de cada controlador, con la cancelación de la
    suscripción al desmontar la pantalla.
    \item \textbf{Manejo de errores de las suscripciones.} Se añadió un
    manejador por defecto que, ante un error de lectura, registra un aviso y
    entrega un resultado vacío para que la interfaz no quede esperando
    indefinidamente. El error de permisos que se produce al cerrar sesión con
    suscripciones activas se reconoce como esperado y se ignora.
    \item \textbf{Limpieza de campos indefinidos.} Firestore rechaza los
    campos con valor \archivo{undefined}; en lugar de exigir que cada
    repositorio lo tenga en cuenta, la limpieza se centralizó en
    \archivo{firestoreService}, el único archivo que usa el SDK.
    \item \textbf{Validación en el controlador.} Todas las operaciones de
    escritura validan sus datos con Zod antes de llegar al repositorio, y las
    reglas de seguridad repiten las validaciones críticas en el servidor.
    \item \textbf{Reglas alineadas con el modelo real.} Las reglas de
    Firestore se revisaron campo por campo contra las entidades y los
    repositorios, de modo que cada escritura que realiza la aplicación
    satisface su regla y ninguna otra es aceptada.
\end{itemize}

\subsection{Métricas del código}

La Tabla~\ref{tab:metricas} resume el tamaño del código fuente por capa al
cierre de la estabilización.

\begin{table}[H]
\centering
\caption{Tamaño del código fuente por capa}
\label{tab:metricas}
\begin{tabular}{|L{3.6cm}|L{4.6cm}|C{1.8cm}|C{1.8cm}|}
\hline
\rowcolor{encabezado} \textbf{Capa} & \textbf{Carpetas} & \textbf{Archivos} & \textbf{Líneas} \\
\hline
Vista & \archivo{app/}, \archivo{components/} & 32 & 4\,860 \\ \hline
Controlador & \archivo{controllers/}, \archivo{utils/}, \archivo{store/},
\archivo{hooks/} & 15 & 1\,399 \\ \hline
Modelo & \archivo{models/}, \archivo{schemas/} & 21 & 1\,165 \\ \hline
Servicios & \archivo{services/}, \archivo{config/} & 6 & 514 \\ \hline
Otros & \archivo{constants/}, \archivo{types/} & 2 & 45 \\
\hline
\textbf{Total} & & \textbf{76} & \textbf{7\,983} \\
\hline
\end{tabular}
\fuente{Conteo de archivos \archivo{.ts} y \archivo{.tsx} de la carpeta
\archivo{src/}, incluidos comentarios y líneas en blanco.}
\end{table}

\subsection{Verificación estática}

Como criterio de cierre de la fase se estableció una secuencia de
verificaciones automáticas que el código debe superar
(Tabla~\ref{tab:verificacion-estatica}).

\begin{table}[H]
\centering
\caption{Verificaciones estáticas del código}
\label{tab:verificacion-estatica}
\begin{tabularx}{\textwidth}{|L{4.3cm}|Y|L{3.2cm}|}
\hline
\rowcolor{encabezado} \textbf{Verificación} & \textbf{Qué detecta} & \textbf{Resultado} \\
\hline
\archivo{tsc --noEmit} & Errores de tipos en todo el proyecto (TypeScript
estricto). & 0 errores \\ \hline
\archivo{eslint .} & Errores y malas prácticas, incluidas las reglas de
\emph{hooks} de React, sin reglas desactivadas. &
\pendiente{registrar} \\ \hline
\archivo{prettier --check} & Formato no uniforme del código. & 75 de 76
archivos conformes \\ \hline
\archivo{expo export -p web} & Rutas duplicadas, importaciones rotas y
dependencias faltantes al empaquetar. & \pendiente{registrar} \\
\hline
\end{tabularx}
\end{table}
